@c -*-texinfo-*-

@c  Copyright (c) 2018 Blake McBride
@c  All rights reserved.
@c
@c  Redistribution and use in source and binary forms, with or without
@c  modification, are permitted provided that the following conditions are
@c  met:
@c
@c  1. Redistributions of source code must retain the above copyright
@c  notice, this list of conditions and the following disclaimer.
@c
@c  2. Redistributions in binary form must reproduce the above copyright
@c  notice, this list of conditions and the following disclaimer in the
@c  documentation and/or other materials provided with the distribution.
@c
@c  THIS SOFTWARE IS PROVIDED BY THE COPYRIGHT HOLDERS AND CONTRIBUTORS
@c  "AS IS" AND ANY EXPRESS OR IMPLIED WARRANTIES, INCLUDING, BUT NOT
@c  LIMITED TO, THE IMPLIED WARRANTIES OF MERCHANTABILITY AND FITNESS FOR
@c  A PARTICULAR PURPOSE ARE DISCLAIMED. IN NO EVENT SHALL THE COPYRIGHT
@c  HOLDER OR CONTRIBUTORS BE LIABLE FOR ANY DIRECT, INDIRECT, INCIDENTAL,
@c  SPECIAL, EXEMPLARY, OR CONSEQUENTIAL DAMAGES (INCLUDING, BUT NOT
@c  LIMITED TO, PROCUREMENT OF SUBSTITUTE GOODS OR SERVICES; LOSS OF USE,
@c  DATA, OR PROFITS; OR BUSINESS INTERRUPTION) HOWEVER CAUSED AND ON ANY
@c  THEORY OF LIABILITY, WHETHER IN CONTRACT, STRICT LIABILITY, OR TORT
@c  (INCLUDING NEGLIGENCE OR OTHERWISE) ARISING IN ANY WAY OUT OF THE USE
@c  OF THIS SOFTWARE, EVEN IF ADVISED OF THE POSSIBILITY OF SUCH DAMAGE.

@node SQL Query Generator

@html
@include style.css
@end html

@chapter SQL Query Generator

In systems with many interrelated SQL tables, writing queries with the
correct joins is complex and error-prone.  Developers must manually
trace foreign key relationships, determine join paths, and construct
multi-table queries --- work that the schema itself already encodes.

@emph{Kiss} provides an automatic SQL query generator that solves this
problem.  The developer specifies what columns to select, what
conditions to apply, and how to order the results.  The system
determines the join path automatically and produces the correct SQL.

All classes are in the @code{org.kissweb.database} package:

@table @code
@item SchemaGraph
Models the database schema as a graph where tables are nodes and
foreign key relationships are edges.  Finds the shortest join path
between tables using BFS (breadth-first search).
@item QueryBuilder
Fluent API to specify SELECT, WHERE, ORDER BY, GROUP BY, HAVING, and
explicit joins.  Uses a @code{SchemaGraph} to automatically determine the
required joins and produce the final SQL.
@item SchemaGraph.Edge
Represents a single or composite foreign key relationship between two
tables.
@end table

@section Quick Start

@subsection Building the Schema Graph

The schema graph must be built once at application startup.  There
are several ways to create it.

@subsubsection From Live Database Metadata

@example
SchemaGraph graph = SchemaGraph.fromDatabase(conn);
@end example

@noindent
This reads JDBC @code{DatabaseMetaData} for all tables and their foreign
keys.  Composite (multi-column) foreign keys are automatically grouped
by constraint name.  Works on all five supported databases (PostgreSQL,
MySQL, SQLite, SQL Server, Oracle).

@subsubsection Programmatic Declaration

@example
SchemaGraph graph = new SchemaGraph();
graph.addForeignKey("employee", "department_id",
                    "department", "department_id");
graph.addForeignKey("order_line",
    new String[]@{"order_id", "product_id"@},
    "order_product",
    new String[]@{"order_id", "product_id"@});
@end example

@subsubsection Hybrid

@example
SchemaGraph graph = SchemaGraph.fromDatabase(conn);
graph.addForeignKey("audit_log", "user_id",
                    "employee", "employee_id");
@end example

@subsubsection Cached (for Fast Startup)

@example
// First run: read from database, save cache
SchemaGraph graph = SchemaGraph.fromDatabase(conn);
graph.saveToFile("schema-cache.txt");

// Subsequent runs: load from cache (sub-millisecond)
SchemaGraph graph = SchemaGraph.loadFromFile("schema-cache.txt");
@end example

@subsection Building and Executing a Query

@example
List<Record> records = new QueryBuilder(graph)
    .select("employee.first_name", "employee.last_name")
    .select("department.name AS dept_name")
    .select("building.address")
    .where("project.project_id = ?", projectId)
    .where("employee.active = 'Y'")
    .orderBy("employee.last_name")
    .orderBy("employee.first_name")
    .fetchAll(conn);

for (Record r : records) @{
    String name = r.getString("first_name");
    String dept = r.getString("dept_name");
@}
@end example

@noindent
The system automatically generates SQL such as:

@example
SELECT employee.first_name, employee.last_name,
       department.name AS dept_name, building.address
FROM project
JOIN project_assignment
  ON project_assignment.project_id = project.project_id
JOIN employee
  ON project_assignment.employee_id = employee.employee_id
JOIN department
  ON employee.department_id = department.department_id
JOIN building
  ON department.building_id = building.building_id
WHERE project.project_id = ?
  AND employee.active = 'Y'
ORDER BY employee.last_name, employee.first_name
@end example

@section SchemaGraph API

@subsection Construction

@deffn Method SchemaGraph ()
Create an empty schema graph for programmatic population.
@end deffn

@deffn {Static Method} SchemaGraph.fromDatabase conn
Build a schema graph by reading all tables and foreign keys from JDBC
@code{DatabaseMetaData}.  Called once at application startup.

@strong{Parameters:}
@itemize @bullet
@item
@code{conn} --- a Kiss @code{Connection} object
@end itemize

@strong{Returns:} a populated @code{SchemaGraph}

@strong{Throws:} @code{SQLException} if database metadata cannot be read
@end deffn

@deffn {Static Method} SchemaGraph.loadFromFile path
Load a schema graph from a cache file previously saved with
@code{saveToFile()}.

@strong{Parameters:}
@itemize @bullet
@item
@code{path} --- path to the cache file
@end itemize

@strong{Returns:} a populated @code{SchemaGraph}

@strong{Throws:} @code{IOException} if the file cannot be read or has an invalid format
@end deffn

@subsection Populating

@deffn Method addTable tableName
Declare a table in the graph.  This is optional --- tables are
automatically created when referenced by @code{addForeignKey}.
@end deffn

@deffn Method addForeignKey fromTable fromColumn toTable toColumn
Declare a single-column foreign key relationship between two tables.
Both tables are automatically added to the graph if not already present.
Duplicate edges are ignored.

@strong{Parameters:}
@itemize @bullet
@item
@code{fromTable} --- the child table containing the FK column
@item
@code{fromColumn} --- the FK column in the child table
@item
@code{toTable} --- the parent/referenced table
@item
@code{toColumn} --- the referenced column (usually the primary key)
@end itemize
@end deffn

@deffn Method addForeignKey fromTable fromColumns[] toTable toColumns[]
Declare a composite (multi-column) foreign key relationship.
Both tables are automatically added to the graph if not already present.
Duplicate edges are ignored.

@strong{Parameters:}
@itemize @bullet
@item
@code{fromTable} --- the child table containing the FK columns
@item
@code{fromColumns} --- array of FK column names in the child table
@item
@code{toTable} --- the parent/referenced table
@item
@code{toColumns} --- array of referenced column names
@end itemize
@end deffn

@subsection Querying the Graph

@deffn Method hasTable tableName
Check whether a table exists in the graph.

@strong{Returns:} @code{true} if the table has been added
@end deffn

@deffn Method getTables
Return all table names in the graph.

@strong{Returns:} an unmodifiable @code{Set<String>} of table names
@end deffn

@deffn Method getEdges tableName
Return all foreign key edges incident on a given table.

@strong{Returns:} unmodifiable list of @code{Edge} objects, or an empty list if the table is not in the graph
@end deffn

@deffn Method findJoinPath tables rootTable
Find the shortest join path connecting all the given tables using BFS.
The algorithm picks the root table as the starting point, then greedily
adds the shortest BFS path from the current tree to each remaining
table (a standard Steiner tree approximation).

@strong{Parameters:}
@itemize @bullet
@item
@code{tables} --- @code{Set<String>} of tables that must be connected
@item
@code{rootTable} --- the preferred root table (used as FROM), or @code{null}
@end itemize

@strong{Returns:} an ordered list of @code{Edge} objects representing the joins

@strong{Throws:} @code{SQLException} if no path exists between some tables
@end deffn

@subsection Schema Caching

@deffn Method saveToFile path
Write the schema graph to a text file for later loading.

The cache file format is line-based text:
@example
# SchemaGraph cache v1
TABLE employee
TABLE department
FK employee department_id department department_id
FK order_line order_id,product_id order_product order_id,product_id
@end example
@end deffn

@deffn {Static Method} SchemaGraph.loadFromFile path
Read a schema graph from a cache file previously saved with @code{saveToFile()}.
Comments (lines starting with @code{#}) and blank lines are allowed.
@end deffn

@subsection Edge Inner Class

@code{SchemaGraph.Edge} represents a foreign key relationship (single or
composite).

@deffn Method getFromTable
Return the child table name.
@end deffn

@deffn Method getToTable
Return the parent/referenced table name.
@end deffn

@deffn Method getFromColumn
Return the first (or only) FK column name.
@end deffn

@deffn Method getToColumn
Return the first (or only) referenced column name.
@end deffn

@deffn Method getFromColumns
Return all FK column names as @code{List<String>}.
@end deffn

@deffn Method getToColumns
Return all referenced column names as @code{List<String>}.
@end deffn

@deffn Method isComposite
Return @code{true} if this is a composite (multi-column) foreign key.
@end deffn

@deffn Method buildOnCondition fromAlias toAlias
Generate the ON condition clause for this edge.  For composite keys,
column pairs are separated with @code{AND}.

@strong{Parameters:}
@itemize @bullet
@item
@code{fromAlias} --- alias for the from table, or @code{null} to use the table name
@item
@code{toAlias} --- alias for the to table, or @code{null} to use the table name
@end itemize

@strong{Returns:} the ON condition string (e.g.@: @code{employee.dept_id = department.dept_id})
@end deffn

@noindent
All table and column names are case-insensitive (stored lowercase
internally).

@section QueryBuilder API

@subsection SELECT

@deffn Method select tableColumn
Add a column to the SELECT list.  Accepts formats such as
@code{"table.column"}, @code{"table.column AS alias"},
@code{"COUNT(table.column)"}, and @code{"SUM(table.column) AS total"}.

@strong{Returns:} @code{this} for chaining
@end deffn

@deffn Method select tableColumns...
Add multiple columns to the SELECT list.

@strong{Returns:} @code{this} for chaining
@end deffn

@deffn Method distinct
Enable @code{SELECT DISTINCT}.

@strong{Returns:} @code{this} for chaining
@end deffn

@subsection Aggregate Helpers

The following methods add aggregate function calls to the SELECT list.
Each has a variant without an alias and a variant with an alias.

@deffn Method selectCount tableColumn
@deffnx Method selectCount tableColumn alias
Add @code{COUNT(table.column)} or @code{COUNT(table.column) AS alias}.
@end deffn

@deffn Method selectSum tableColumn
@deffnx Method selectSum tableColumn alias
Add @code{SUM(table.column)} or @code{SUM(table.column) AS alias}.
@end deffn

@deffn Method selectAvg tableColumn
@deffnx Method selectAvg tableColumn alias
Add @code{AVG(table.column)} or @code{AVG(table.column) AS alias}.
@end deffn

@deffn Method selectMin tableColumn
@deffnx Method selectMin tableColumn alias
Add @code{MIN(table.column)} or @code{MIN(table.column) AS alias}.
@end deffn

@deffn Method selectMax tableColumn
@deffnx Method selectMax tableColumn alias
Add @code{MAX(table.column)} or @code{MAX(table.column) AS alias}.
@end deffn

@subsection WHERE

@deffn Method where condition params...
Add a WHERE condition.  Multiple calls at the same level are combined
with the current group's operator (AND by default).
Use @code{?} for parameter placeholders.

@strong{Example:}
@example
.where("employee.active = ?", "Y")
.where("employee.salary > ?", 50000)
@end example
@end deffn

@subsection OR / AND Grouping

By default, multiple @code{where()} calls are combined with AND.
Use @code{startOr()} and @code{endOr()} to group conditions with OR:

@example
String sql = new QueryBuilder(graph)
    .select("employee.first_name", "employee.last_name")
    .startOr()
        .where("employee.department_id = ?", 10)
        .where("employee.department_id = ?", 20)
    .endOr()
    .where("employee.active = 'Y'")
    .build();
// WHERE (department_id = ? OR department_id = ?)
//   AND employee.active = 'Y'
@end example

@noindent
OR groups can be nested with explicit AND groups:

@example
String sql = new QueryBuilder(graph)
    .select("employee.first_name")
    .startOr()
        .where("employee.salary > ?", 100000)
        .startAnd()
            .where("employee.department_id = ?", 5)
            .where("employee.level = ?", "senior")
        .endAnd()
    .endOr()
    .build();
// WHERE (salary > ? OR (department_id = ? AND level = ?))
@end example

@noindent
Conditions can be added conditionally (the most common pattern in
real codebases):

@example
QueryBuilder qb = new QueryBuilder(graph)
    .select("employee.first_name", "employee.last_name");
qb.startOr();
if (dept1 != null) qb.where("employee.department_id = ?", dept1);
if (dept2 != null) qb.where("employee.department_id = ?", dept2);
qb.endOr();
@end example

@noindent
Empty OR/AND groups (no conditions inside) are silently omitted from
the SQL.

@deffn Method startOr
Begin an OR group --- conditions added between @code{startOr()} and
@code{endOr()} are combined with OR.
@end deffn

@deffn Method endOr
End the current OR group and return to the parent group.
@end deffn

@deffn Method startAnd
Begin an AND group.  Useful inside an OR group to explicitly group
conditions with AND.
@end deffn

@deffn Method endAnd
End the current AND group and return to the parent group.
@end deffn

@subsection Subqueries in WHERE

@subsubsection IN / NOT IN with Subquery

Use @code{whereIn()} and @code{whereNotIn()} with a nested
@code{QueryBuilder} or raw SQL:

@example
QueryBuilder sub = new QueryBuilder(graph)
    .select("project_assignment.employee_id")
    .where("project_assignment.project_id = ?", projectId);

String sql = new QueryBuilder(graph)
    .select("employee.first_name", "employee.last_name")
    .whereIn("employee.employee_id", sub)
    .build();
// WHERE employee.employee_id IN
//   (SELECT project_assignment.employee_id ...)
@end example

@noindent
With raw SQL:

@example
String sql = new QueryBuilder(graph)
    .select("employee.first_name")
    .whereIn("employee.employee_id",
        "SELECT person_id FROM timesheet WHERE work_date = ?",
        20260101)
    .build();
@end example

@subsubsection EXISTS / NOT EXISTS

@example
QueryBuilder sub = new QueryBuilder(graph)
    .select("1")
    .where("project_assignment.employee_id = employee.employee_id")
    .where("project_assignment.project_id = ?", 42);

String sql = new QueryBuilder(graph)
    .select("employee.first_name")
    .whereExists(sub)
    .build();
// WHERE EXISTS (SELECT 1 FROM project_assignment WHERE ...)
@end example

@noindent
With raw SQL:

@example
String sql = new QueryBuilder(graph)
    .select("department.name")
    .whereExists("SELECT 1 FROM employee "
       + "WHERE employee.department_id = department.department_id "
       + "AND employee.active = 'Y'")
    .build();
@end example

@subsubsection Subquery Methods

@deffn Method whereIn column subquery
Add @code{column IN (SELECT ...)} using a @code{QueryBuilder} subquery.
@end deffn

@deffn Method whereIn column rawSQL params...
Add @code{column IN (raw SQL)} with parameter values.
@end deffn

@deffn Method whereNotIn column subquery
Add @code{column NOT IN (SELECT ...)} using a @code{QueryBuilder} subquery.
@end deffn

@deffn Method whereNotIn column rawSQL params...
Add @code{column NOT IN (raw SQL)} with parameter values.
@end deffn

@deffn Method whereExists subquery
Add @code{EXISTS (SELECT ...)} using a @code{QueryBuilder} subquery.
@end deffn

@deffn Method whereExists rawSQL params...
Add @code{EXISTS (raw SQL)} with parameter values.
@end deffn

@deffn Method whereNotExists subquery
Add @code{NOT EXISTS (SELECT ...)} using a @code{QueryBuilder} subquery.
@end deffn

@deffn Method whereNotExists rawSQL params...
Add @code{NOT EXISTS (raw SQL)} with parameter values.
@end deffn

@noindent
Subquery parameters are merged into the outer query's parameter list
in the correct positional order.

@subsection ORDER BY

@deffn Method orderBy tableColumn
Add an ascending ORDER BY column.
@end deffn

@deffn Method orderByDesc tableColumn
Add a descending ORDER BY column.
@end deffn

@subsection GROUP BY / HAVING

@deffn Method groupBy tableColumn
Add a GROUP BY column.
@end deffn

@deffn Method having condition params...
Add a HAVING condition.  Multiple calls are combined with @code{AND}.
@end deffn

@subsection Explicit Joins

Explicit joins give the developer control over exactly which columns
are used for a join.  They are necessary for self-joins, and useful
when multiple foreign key paths exist between two tables and the
developer needs to select a specific one.

@subsubsection INNER JOIN

@deffn Method join fromTable fromColumn toTable toColumn
Single-column INNER JOIN without alias.
@end deffn

@deffn Method join fromTable fromColumn toTable toColumn alias
Single-column INNER JOIN with alias.  The alias is required for self-joins.
@end deffn

@deffn Method join fromTable fromColumns[] toTable toColumns[] alias
Composite (multi-column) INNER JOIN, optionally with alias.
@end deffn

@subsubsection LEFT JOIN

@deffn Method leftJoin fromTable fromColumn toTable toColumn
Single-column LEFT JOIN without alias.
@end deffn

@deffn Method leftJoin fromTable fromColumn toTable toColumn alias
Single-column LEFT JOIN with alias.
@end deffn

@deffn Method leftJoin fromTable fromColumns[] toTable toColumns[] alias
Composite LEFT JOIN, optionally with alias.
@end deffn

@subsubsection RIGHT JOIN

@deffn Method rightJoin fromTable fromColumn toTable toColumn
Single-column RIGHT JOIN without alias.
@end deffn

@deffn Method rightJoin fromTable fromColumn toTable toColumn alias
Single-column RIGHT JOIN with alias.
@end deffn

@deffn Method rightJoin fromTable fromColumns[] toTable toColumns[] alias
Composite RIGHT JOIN, optionally with alias.
@end deffn

@subsection LIMIT

@deffn Method limit max
Limit the number of rows returned.  Uses @code{Connection.limit()} to
produce database-specific SQL.
@end deffn

@subsection Common Table Expressions (CTEs)

Use @code{with()} to prepend CTE blocks to the query:

@example
String sql = new QueryBuilder(graph)
    .with("last_timesheet",
          "SELECT person_id, MAX(work_date) work_date "
        + "FROM timesheet WHERE billable = 'Y' "
        + "GROUP BY person_id")
    .select("employee.first_name", "employee.last_name")
    .select("last_timesheet.work_date")
    .leftJoin("employee", "employee_id",
              "last_timesheet", "person_id")
    .build();
// WITH last_timesheet AS (SELECT ...) SELECT ...
@end example

@noindent
Multiple CTEs are supported:

@example
String sql = new QueryBuilder(graph)
    .with("cte1", "SELECT employee_id FROM employee "
        + "WHERE active = 'Y'")
    .with("cte2", "SELECT employee_id, MAX(work_date) "
        + "last_date FROM timesheet GROUP BY employee_id")
    .select("cte1.employee_id")
    .select("cte2.last_date")
    .leftJoin("cte1", "employee_id",
              "cte2", "employee_id")
    .build();
// WITH cte1 AS (...), cte2 AS (...) SELECT ...
@end example

@noindent
A CTE can also use a @code{QueryBuilder} as its body:

@example
QueryBuilder cteSub = new QueryBuilder(graph)
    .select("employee.employee_id")
    .where("employee.active = ?", "Y");

String sql = new QueryBuilder(graph)
    .with("active_emps", cteSub)
    .select("active_emps.employee_id")
    .build();
@end example

@deffn Method with name rawSQL params...
Add a CTE with a raw SQL body and optional parameters.
@end deffn

@deffn Method with name subquery
Add a CTE whose body is built from a @code{QueryBuilder}.
@end deffn

@noindent
CTE names are treated like aliases --- they are excluded from the
schema graph path finding.  CTE parameters appear before WHERE
parameters in the final parameter list.

@subsection UNION

Combine multiple queries with @code{union()} or @code{unionAll()}:

@example
QueryBuilder q1 = new QueryBuilder(graph)
    .select("employee.first_name", "employee.last_name")
    .select("'Employee' AS person_type")
    .where("employee.active = 'Y'");

QueryBuilder q2 = new QueryBuilder(graph)
    .select("employee.first_name", "employee.last_name")
    .select("'Contractor' AS person_type")
    .where("employee.active = 'N'");

String sql = q1.unionAll(q2)
    .orderBy("last_name")
    .build();
// SELECT ... UNION ALL SELECT ... ORDER BY last_name
@end example

@deffn Method union other
Append a UNION (removes duplicates) with another @code{QueryBuilder}.
@end deffn

@deffn Method unionAll other
Append a UNION ALL (keeps duplicates) with another @code{QueryBuilder}.
@end deffn

@noindent
ORDER BY and LIMIT on a UNION apply to the combined result.  Multiple
unions can be chained.  Parameters from all parts are merged in the
correct positional order.

@subsection Build and Execute

@deffn Method build
Generate the SQL query string.  After calling this method,
@code{getParameters()} returns the ordered parameter values.

@strong{Returns:} the generated SQL string

@strong{Throws:} @code{SQLException} if no join path exists between the referenced tables, or if no SELECT columns were specified
@end deffn

@deffn Method getParameters
Return the ordered list of @code{?} placeholder values.  Must be called
after @code{build()}.

@strong{Returns:} unmodifiable @code{List<Object>}

@strong{Throws:} @code{IllegalStateException} if called before @code{build()}
@end deffn

@deffn Method fetchAll conn
Build and execute the query, returning all results as @code{Record}
objects.

@strong{Returns:} @code{List<Record>}

@strong{Throws:} @code{Exception} on database or query building errors
@end deffn

@deffn Method fetchOne conn
Build and execute the query, returning the first result or @code{null}.

@strong{Returns:} @code{Record} or @code{null}

@strong{Throws:} @code{Exception} on database or query building errors
@end deffn

@deffn Method fetchAllJSON conn
Build and execute the query, returning results as a @code{JSONArray}.

@strong{Returns:} @code{JSONArray}

@strong{Throws:} @code{Exception} on database or query building errors
@end deffn

@section Examples

@subsection Basic Multi-Table Query

@example
String sql = new QueryBuilder(graph)
    .select("employee.first_name", "employee.last_name")
    .select("department.name")
    .select("building.address")
    .where("employee.active = ?", "Y")
    .orderBy("employee.last_name")
    .build();
@end example

@noindent
The builder automatically determines the join path from
@code{employee} through @code{department} to @code{building}
based on the foreign key relationships in the schema graph.

@subsection Self-Join (Employee to Manager)

@example
String sql = new QueryBuilder(graph)
    .select("employee.first_name")
    .select("mgr.first_name AS manager_name")
    .join("employee", "manager_id",
          "employee", "employee_id", "mgr")
    .build();
// ... JOIN employee mgr
//       ON employee.manager_id = mgr.employee_id
@end example

@subsection Multiple Self-Joins

@example
String sql = new QueryBuilder(graph)
    .select("employee.first_name")
    .select("mgr.first_name AS manager")
    .select("mentor.first_name AS mentor_name")
    .join("employee", "manager_id",
          "employee", "employee_id", "mgr")
    .join("employee", "mentor_id",
          "employee", "employee_id", "mentor")
    .build();
@end example

@subsection LEFT JOIN with Aggregate

@example
String sql = new QueryBuilder(graph)
    .select("department.name")
    .selectCount("employee.employee_id", "emp_count")
    .leftJoin("department", "department_id",
              "employee", "department_id")
    .groupBy("department.name")
    .build();
// SELECT department.name,
//        COUNT(employee.employee_id) AS emp_count
// FROM department
// LEFT JOIN employee
//   ON department.department_id = employee.department_id
// GROUP BY department.name
@end example

@subsection Multiple Aggregates

@example
String sql = new QueryBuilder(graph)
    .select("department.name")
    .selectCount("employee.employee_id", "emp_count")
    .selectAvg("employee.salary", "avg_salary")
    .selectSum("employee.salary", "total_salary")
    .selectMin("employee.salary", "min_salary")
    .selectMax("employee.salary", "max_salary")
    .groupBy("department.name")
    .having("COUNT(employee.employee_id) > ?", 5)
    .build();
@end example

@subsection Composite Foreign Key Join

@example
String sql = new QueryBuilder(graph)
    .select("order_line.qty")
    .select("pricing.price")
    .join("order_line",
          new String[]@{"product_id", "region_id"@},
          "pricing",
          new String[]@{"product_id", "region_id"@}, null)
    .build();
// ... JOIN pricing
//   ON order_line.product_id = pricing.product_id
//  AND order_line.region_id = pricing.region_id
@end example

@subsection Picking a Specific Foreign Key Path

When a table has multiple foreign keys pointing to the same table,
use an explicit join to select which one:

@example
// employee has both department_id and home_department_id
String sql = new QueryBuilder(graph)
    .select("employee.name")
    .select("department.name AS home_dept")
    .join("employee", "home_department_id",
          "department", "department_id")
    .build();
@end example

@subsection Schema Caching for Fast Startup

@example
// First run: read from database, save cache
SchemaGraph graph = SchemaGraph.fromDatabase(conn);
graph.saveToFile("schema-cache.txt");

// Subsequent runs: load from cache (sub-millisecond)
SchemaGraph graph = SchemaGraph.loadFromFile("schema-cache.txt");

// Hybrid: load cache, then add extra relationships
SchemaGraph graph = SchemaGraph.loadFromFile("schema-cache.txt");
graph.addForeignKey("audit_log", "user_id",
                    "employee", "employee_id");
@end example

@subsection OR Conditions with Dynamic Building

@example
QueryBuilder qb = new QueryBuilder(graph)
    .select("employee.first_name", "employee.last_name");
qb.startOr();
if (dept1 != null)
    qb.where("employee.department_id = ?", dept1);
if (dept2 != null)
    qb.where("employee.department_id = ?", dept2);
qb.endOr();
qb.where("employee.active = 'Y'");
String sql = qb.build();
// WHERE (department_id = ? OR department_id = ?)
//   AND employee.active = 'Y'
@end example

@subsection Subquery in WHERE

@example
QueryBuilder sub = new QueryBuilder(graph)
    .select("project_assignment.employee_id")
    .where("project_assignment.project_id = ?", projId);

String sql = new QueryBuilder(graph)
    .select("employee.first_name")
    .where("employee.active = 'Y'")
    .whereIn("employee.employee_id", sub)
    .build();
// WHERE employee.active = 'Y'
//   AND employee.employee_id IN
//     (SELECT project_assignment.employee_id ...)
@end example

@subsection CTE with LEFT JOIN

@example
String sql = new QueryBuilder(graph)
    .with("last_ts",
          "SELECT person_id, MAX(work_date) work_date "
        + "FROM timesheet WHERE billable = 'Y' "
        + "GROUP BY person_id")
    .select("employee.first_name")
    .select("last_ts.work_date")
    .leftJoin("employee", "employee_id",
              "last_ts", "person_id")
    .build();
@end example

@subsection UNION ALL with ORDER BY

@example
QueryBuilder q1 = new QueryBuilder(graph)
    .select("employee.first_name")
    .select("'Active' AS status")
    .where("employee.active = 'Y'");

QueryBuilder q2 = new QueryBuilder(graph)
    .select("employee.first_name")
    .select("'Inactive' AS status")
    .where("employee.active = 'N'");

String sql = q1.unionAll(q2)
    .orderBy("first_name")
    .build();
@end example

@section How It Works

@subsection Join Path Finding (BFS)

When @code{build()} is called, the query builder:

@enumerate
@item
Collects all distinct table names mentioned in SELECT, WHERE, ORDER BY,
GROUP BY, and HAVING clauses.
@item
Picks the root table (first WHERE table, or first SELECT table).
@item
For each remaining table, runs BFS from the set of already-connected
tables to find the shortest path through the schema graph.
@item
Merges overlapping paths (deduplicates edges).  The result is a tree
of joins connecting all required tables.
@item
If any table is unreachable, throws @code{SQLException} naming the
disconnected tables.
@end enumerate

@noindent
BFS is optimal because all foreign key edges have equal weight.  It
runs in O(V+E) time, which is sub-millisecond even for 300+ tables.

@subsection Explicit Join Isolation

Tables that are the target (@code{toTable}) of an explicit @code{join()},
@code{leftJoin()}, or @code{rightJoin()} are excluded from the automatic
join graph traversal.  This means:

@itemize @bullet
@item
Explicit joins work even for tables not in the schema graph.
@item
Self-joins are handled correctly (same table as both from and to, with
an alias).
@item
The developer controls exactly how explicit-join tables are reached.
@end itemize

@subsection Join Type Rules

@itemize @bullet
@item
Auto-discovered joins (from schema graph path finding) are always @code{INNER JOIN}.
@item
Explicit joins use the type specified: @code{join()} produces @code{JOIN},
@code{leftJoin()} produces @code{LEFT JOIN}, and @code{rightJoin()} produces @code{RIGHT JOIN}.
@end itemize

@section Edge Cases and Limitations

@table @asis
@item Tables with no FK path between them
@code{build()} throws @code{SQLException} naming the disconnected tables.
@item Composite foreign keys (multi-column)
Fully supported.  @code{Edge} holds column lists; the ON condition emits
@code{AND}-separated column pairs.
@item Self-joins
Supported via explicit @code{join()} with an alias.
@item Multiple FK paths between two tables
Auto-join picks the first path found.  Use an explicit @code{join()} to
select a specific path.
@item Tables only reachable via explicit join
Works correctly --- explicit join target tables bypass graph path finding.
@item Schema-qualified table names
Supported --- table names include the schema prefix throughout.
@item Aggregate queries without GROUP BY
The developer's responsibility to ensure correct SQL semantics.
@item OR conditions
Supported via @code{startOr()} / @code{endOr()} grouping with nesting.
@item Subqueries (IN, EXISTS)
Supported via @code{whereIn()}, @code{whereExists()}, etc.
@item CTEs (WITH ... AS)
Supported via @code{with()}.  CTE names are excluded from auto-join path finding.
@item UNION / UNION ALL
Supported via @code{union()} / @code{unionAll()}.
@item INTERSECT, EXCEPT
Not supported; use raw SQL.
@item FULL OUTER JOIN
Not supported.
@end table

@section Performance

@itemize @bullet
@item
@strong{Schema graph construction from DB:} One-time cost, typically 1--3
seconds for 300 tables.
@item
@strong{Schema graph load from cache file:} Sub-millisecond.
@item
@strong{Path finding:} Sub-millisecond BFS on 300 nodes.
@item
@strong{Memory:} Adjacency list for 300 tables with ~500 FK relationships
uses less than 1 MB.
@item
@strong{Thread safety:} @code{SchemaGraph} is safe for concurrent read
operations after construction.
@end itemize

@section Source Files

@table @code
@item src/main/core/org/kissweb/database/SchemaGraph.java
Schema graph model: table/edge management, BFS path finding, schema
caching.
@item src/main/core/org/kissweb/database/QueryBuilder.java
Fluent query builder: SELECT, WHERE, ORDER BY, GROUP BY, HAVING,
explicit joins, aggregate helpers, build, and execute methods.
@item src/test/core/org/kissweb/database/QueryBuilderTest.java
108 unit tests covering all functionality.
@end table
